\newpage
\section{交点}
\label{sec:intersections}

函数 \tkzFct{tkz-elements}{intersection} 是一个必不可少的几何工具。它计算两个几何对象的交点，这两个对象可以属于以下类：

\begin{itemize}
  \item \tkzClass{line} 与 \tkzClass{line}，
  \item \tkzClass{line} 与 \tkzClass{circle}，
  \item \tkzClass{circle} 与 \tkzClass{circle}，
  \item \tkzClass{line} 与 \tkzClass{conic}。
\end{itemize}

注意，\tkzClass{circle} 是与 \tkzClass{conic} 不同的类。
\code{conic} 类包括抛物线、双曲线和椭圆。

\medskip

该函数接收一对对象作为参数，与顺序无关。结果通常由一个或两个点组成，取决于几何构型。若没有交点，函数不返回任何值（\code{nil}）；结果可用 \code{if not x then} 来判断。任何其他对象对（例如一个圆和一条圆锥曲线，或两条圆锥曲线）都会引发错误。

\medskip

当存在两个交点，而你又：
\begin{itemize}
  \item 已经知道其中一个，或
  \item 想选取距某个参考点最近的那一个，
\end{itemize}
时，可以在 \code{opts} 表中使用可选键。

此外，数值稳健性可通过一个可选的容差参数 \code{eps} 控制，它同样存放在 \code{opts} 中。

\subsection{可选参数：\code{known}、\code{near} 和 \code{eps}}

函数 \tkzFct{tkz-elements}{intersection(X, Y, opts)} 计算两个几何对象 \code{X} 和 \code{Y} 之间的交点。可选表 \code{opts} 可包含以下键：

\begin{itemize}
  \item \tkzKey{known} \tkzOptArg{point} \\
  当两个交点之一已经已知时使用本选项。
  函数会把\emph{另一个}解作为第一个返回值。
  这在其中一个交点已经确定的作图中特别有用：

  \begin{mybox}
  \begin{verbatim}
  z.X, z.Y = intersection(C1, C2, {known = z.Y})
  \end{verbatim}
  \end{mybox}

  若 \code{z.Y} 是交点之一，函数保证 \code{z.X} 得到另一个。

  \item \tkzKey{near} \tkzOptArg{point} \\
  使用本选项时，函数先返回距给定参考点最近的那个解：

  \begin{mybox}
  \begin{verbatim}
  z.A, z.B = intersection(C1, L, {near = z.O})
  \end{verbatim}
  \end{mybox}

  这里，\code{z.A} 将是距 \code{z.O} 最近的交点。

  \item \tkzKey{eps} \tkzOptArg{number} \\
  这个可选键设定内部用于判断对象是相交、相触还是相切时的数值容差。
  若省略，则使用全局值 \code{tkz.epsilon}
  （接受 \code{EPS} 作为别名）：

  \begin{mybox}
  \begin{verbatim}
  z.A, z.B = intersection(C1, C2, {eps = 1e-6})
  \end{verbatim}
  \end{mybox}

  本选项在涉及近似相切或微小数值不确定性的构型中可提高稳健性。

\end{itemize}

若未提供这些选项中的任何一个，交点将按几何计算所确定的默认顺序返回，并使用全局容差 \code{tkz.epsilon}。

\medskip

\texttt{注意： } 这些选项只影响返回多个点的情形。对于单点相交（例如相切），结果不受影响。不过，\code{eps} 选项始终控制内部使用的数值容差。

\medskip
\texttt{兼容性说明： }
为方便起见并向后兼容，\code{intersection} 的第三个参数也可以是一个数值容差：
\begin{mybox}
\begin{verbatim}
z.A, z.B = intersection(C1, C2, 1e-5)
\end{verbatim}
\end{mybox}
此时，该数被解释为 \code{eps} 的值，且不提供其他选项。
推荐的现代写法是：
\begin{mybox}
\begin{verbatim}
z.A, z.B = intersection(C1, C2, {eps = 1e-5})
\end{verbatim}
\end{mybox}




\subsection{直线与直线}
\label{sub:line_line}

结果是一个点；若两直线平行，则什么也不返回（|nil|）。

\begin{minipage}{0.6\textwidth}
\begin{verbatim}
\directlua{
 init_elements()
 z.A = point(1, -1)
 z.B = point(4, 1)
 z.C = point(2, 1)
 z.D = point(4, -2)
 z.I = point(0, 0)
 L.AB = line(z.A, z.B)
 L.CD = line(z.C, z.D)
 x = intersection(L.AB, L.CD)
 if not x then
 	tex.print("error")
 else
 	z.I = x
end}

\begin{tikzpicture}
   \tkzGetNodes
   \tkzDrawSegments(A,B C,D)
   \tkzDrawPoints(A,B,C,D,I)
   \tkzLabelPoints(A,B,C,D,I)
\end{tikzpicture}
\end{verbatim}
\end{minipage}
\begin{minipage}{0.4\textwidth}
\directlua{
 init_elements()
 z.A = point(1, -1)
 z.B = point(4, 1)
 z.C = point(2, 1)
 z.D = point(4, -2)
 z.I = point(0, 0)
 L.AB = line(z.A, z.B)
 L.CD = line(z.C, z.D)
 x = intersection(L.AB, L.CD)
 if not x then
 	tex.print("error")
 else
 	z.I = x
end}

\begin{tikzpicture}[scale = 2]
   \tkzGetNodes
   \tkzDrawSegments(A,B C,D)
   \tkzDrawPoints(A,B,C,D,I)
   \tkzLabelPoints(A,B,C,D,I)
\end{tikzpicture}
\end{minipage}

\newpage

\subsection{直线与圆}

结果形如：|point,point|；若无交点则什么也不返回（|nil|）。若直线与圆相切，则两个点相同。可以用下划线忽略其中一个点：|_, point| 或 |point, _|。当交给出两个解时，点的顺序由 |(z.p - z.c)| 的辐角决定，其中 |c| 是圆心，|p| 是交点。第一个解对应最小的辐角（辐角在 0 到 $2\pi$ 之间）。

\begin{minipage}{0.6\textwidth}
\begin{verbatim}
\directlua{
 init_elements()
 z.A = point(1, -1)
 z.B = point(1, 2)
 L.AB = line(z.A, z.B)
 z.O = point(2, 1)
 z.D = point(3, 1)
 z.E = point(3, 2)
 L.AE = line(z.A, z.E)
 C.OD = circle(z.O, z.D)
 z.I, _ = intersection(L.AB, C.OD)
 _, z.K = intersection(C.OD, L.AE)}

\begin{tikzpicture}
\tkzGetNodes
   \tkzDrawLines(A,B A,E)
   \tkzDrawCircle(O,D)
   \tkzDrawPoints(A,B,O,D,I,K)
   \tkzLabelPoints[left](A,B,O,D,I,K)
\end{tikzpicture}
\end{verbatim}
\end{minipage}
\begin{minipage}{0.4\textwidth}
\directlua{
 init_elements()
 z.A = point(1, -1)
 z.B = point(1, 2)
 L.AB = line(z.A, z.B)
 z.O = point(2, 1)
 z.D = point(3, 1)
 z.E = point(3, 2)
 L.AE = line(z.A, z.E)
 C.OD = circle(z.O, z.D)
 z.I, _ = intersection(L.AB, C.OD)
 _, z.K = intersection(C.OD, L.AE)}

\begin{tikzpicture}[ scale = 2]
\tkzGetNodes
\tkzDrawLines[add=.1 and .1](A,B A,E)
\tkzDrawCircle(O,D)
\tkzDrawPoints(A,B,O,D,I,K)
\tkzLabelPoints[left](A,B,O,D,I,K)
\end{tikzpicture}
\hfill
\end{minipage}

\newpage

\subsection{圆与圆}

结果形如：|point,point|；若无交点则什么也不返回（|nil|）。若两圆相切，则两个点相同。可以用下划线忽略其中一个点：|_ , point| 或 |point , _|。与直线和圆的交点一样，考察 |z.p-z.c| 的辐角，其中 |c| 是第一个圆的圆心，|p| 是交点。第一个解对应最小的辐角（辐角在 0 到 $2\pi$ 之间）。

\begin{minipage}{0.5\textwidth}
\begin{verbatim}
\directlua{
 init_elements()
 z.A = point(1, 1)
 z.B = point(2, 2)
 z.C = point(3, 3)
 z.D = point(3, 0)
 C.AB = circle(z.A, z.B)
 C.CB = circle(z.C, z.B)
 z.I, _ = intersection(C.AB, C.CB)
 C.DC = circle(z.D, z.C)
 z.J, z.K = intersection(C.DC, C.CB)}
\begin{tikzpicture}
   \tkzGetNodes
   \tkzDrawCircles(A,B C,B D,C)
   \tkzDrawPoints(A,I,C,D,J,K)
   \tkzLabelPoints(A,I,C,D,J,K)
\end{tikzpicture}
\end{verbatim}
\end{minipage}
\begin{minipage}{0.5\textwidth}
\directlua{
 init_elements()
 z.A = point(1, 1)
 z.B = point(2, 2)
 z.C = point(3, 3)
 z.D = point(3, 0)
 C.AB = circle(z.A, z.B)
 C.CB = circle(z.C, z.B)
 z.I, _ = intersection(C.AB, C.CB)
 C.DC = circle(z.D, z.C)
 z.J, z.K = intersection(C.DC, C.CB)}

\begin{tikzpicture}
 \tkzGetNodes
 \tkzDrawCircles(A,B C,B D,C)
 \tkzDrawPoints(A,I,C,D,J,K)
 \tkzLabelPoints(A,I,C,D,J,K)
\end{tikzpicture}
\end{minipage}

另一个示例：\ref{sub:the_figure_pappus_circle}
\newpage

\subsection{直线与圆锥曲线}
\label{sub:line_conic}
下面这个示例比较复杂，但它展示了 Lua 的能力。交点的指定比前一个稍微复杂一些，因为必须考虑表征长轴的那个辐角。原理相同，但要减去这个辐角。具体而言，需要考察椭圆中心与各交点所连直线的斜率，以及长轴的斜率。

对于抛物线或双曲线，两个点按到直线第一个点（|pa|）的距离递增排序。一条平行于抛物线轴、或平行于双曲线渐近线的直线，只与圆锥曲线交于一个点，该点单独返回；渐近线与双曲线不相交。
\vspace{1em}

\begin{minipage}{0.5\textwidth}
\begin{verbatim}
\directlua{
 init_elements()
 z.a = point(5, 2)
 z.b = point(-4, 0)
 L.ab = line(z.a, z.b)
 z.c = L.ab.mid
 z.v = L.ab:point(-0.2)
 local a = tkz.length(z.c, z.v)
 local c = 0.5 * tkz.length(z.a, z.b)
 local e = c / a
 z.K = L.ab:report(a ^ 2 / c, z.c)
 z.Kp = (z.K - z.a):orthogonal(2):at(z.K)
 L.dir = line(z.K, z.Kp)
 CO.EL = conic(z.b, L.dir, e)
 PA.EL = CO.EL:points(0, 1, 50)
 z.m = point(2, 4)
 z.n = point(4, 4)
 L.mn = line(z.m, z.n)
 z.r, z.s = intersection(CO.EL, L.mn)}
 \end{verbatim}
 \end{minipage}
\begin{minipage}{0.5\textwidth}
\directlua{
 init_elements()
 z.a = point(5, 2)
 z.b = point(-4, 0)
 L.ab = line(z.a, z.b)
 z.c = L.ab.mid
 z.v = L.ab:point(-0.2)
 local a = tkz.length(z.c, z.v)
 local c = 0.5 * tkz.length(z.a, z.b)
 local e = c / a
 z.K = L.ab:report(a ^ 2 / c, z.c)
 z.Kp = (z.K - z.a):orthogonal(2):at(z.K)
 L.dir = line(z.K, z.Kp)
 CO.EL = conic(z.b, L.dir, e)
 PA.EL = CO.EL:points(0, 1, 50)
 z.m = point(2, 4)
 z.n = point(4, 4)
 L.mn = line(z.m, z.n)
 z.r, z.s = intersection(CO.EL, L.mn)}

\begin{center}
\begin{tikzpicture}[scale =.5]
   \tkzGetNodes
   \tkzDrawLines[red](a,b r,s)
   \tkzDrawSegments(c,r c,s)
   \tkzDrawPoints(a,b,c,r,s)
   \tkzLabelPoints(a,b,c,r,s)
   \tkzDrawCoordinates[smooth,red](PA.EL)
   \tkzFillAngles[green!30,opacity=.4](v,c,s)
   \tkzFillAngles[green!80,opacity=.4](v,c,r)
\end{tikzpicture}
\end{center}
\end{minipage}

\begin{verbatim}
  \begin{tikzpicture}[scale =.5]
     \tkzGetNodes
     \tkzDrawLines[red](a,b r,s)
     \tkzDrawSegments(c,r c,s)
     \tkzDrawPoints(a,b,c,r,s)
     \tkzLabelPoints(a,b,c,r,s)
     \tkzDrawCoordinates[smooth,red](PA.EL)
     \tkzFillAngles[green!30,opacity=.4](v,c,s)
     \tkzFillAngles[green!80,opacity=.4](v,c,r)
  \end{tikzpicture}
\end{verbatim}


\subsubsection{与圆锥曲线各子类的交点}

\begin{verbatim}
\directlua{
 init_elements()
 z.A = point(0, 0)
 z.B = point(4, -2)
 L.dir = line(z.A, z.B)
 z.F = point(2, 2)
 CO.EL = conic(z.F, L.dir, 0.8)
 CO.PA = conic(z.F, L.dir, 1)
 CO.HY = conic(z.F, L.dir, 1.2)
 PA.EL = CO.EL:points(0, 1, 50)
 PA.PA = CO.PA:points(-5, 5, 50)
 PA.HY = CO.HY:points(-5, 5, 50)
 z.K = CO.EL.K
 z.u, z.v = CO.EL.major_axis:get()
 z.x = L.dir:report(-4, z.K)
 z.y = L.dir:report(4, z.K)
 z.r = point(0, 4)
 z.s = point(4, 1)
 L.rs = line(z.r, z.s)
 z.u_1, z.u_2 = intersection(L.rs, CO.EL)
 z.v_1, z.v_2 = intersection(L.rs, CO.PA)
 z.w_1, z.w_2 = intersection(L.rs, CO.HY)}
\begin{tikzpicture}
 \tkzGetNodes
 \tkzDrawCoordinates[smooth](PA.EL)
 \tkzDrawCoordinates[smooth](PA.PA)
 \tkzDrawCoordinates[smooth](PA.HY)
 \tkzDrawLines[add =.5 and .5](r,s u,v)
 \tkzDrawLines(x,y)
 \tkzDrawPoints[red](u_1,u_2,v_2,v_1,w_1,w_2)
\end{tikzpicture}
\end{verbatim}
\directlua{
 init_elements()
 z.A = point(0, 0)
 z.B = point(4, -2)
 L.dir = line(z.A, z.B)
 z.F = point(2, 2)
 CO.EL = conic(z.F, L.dir, 0.8)
 CO.PA = conic(z.F, L.dir, 1)
 CO.HY = conic(z.F, L.dir, 1.2)
 PA.EL = CO.EL:points(0, 1, 50)
 PA.PA = CO.PA:points(-5, 5, 50)
 PA.HY = CO.HY:points(-5, 5, 50)
 z.K = CO.EL.K
 z.u, z.v = CO.EL.major_axis:get()
 z.x = L.dir:report(-4, z.K)
 z.y = L.dir:report(4, z.K)
 z.r = point(0, 4)
 z.s = point(4, 1)
 L.rs = line(z.r, z.s)
 z.u_1, z.u_2 = intersection(L.rs, CO.EL)
 z.v_1, z.v_2 = intersection(L.rs, CO.PA)
 z.w_1, z.w_2 = intersection(L.rs, CO.HY)}

\begin{center}
  \begin{tikzpicture}[scale =.75]
  \tkzGetNodes
  \tkzDrawCoordinates[smooth](PA.EL)
  \tkzDrawCoordinates[smooth](PA.PA)
  \tkzDrawCoordinates[smooth](PA.HY)
  \tkzDrawLines[add =.5 and .5](r,s u,v)
  \tkzDrawLines(x,y)
  \tkzDrawPoints[red](u_1,u_2,v_2,v_1,w_1,w_2)
  \end{tikzpicture}
\end{center}

\subsubsection{直线与抛物线的交点（详解）}
\label{ssub:intersection_line_parabola_explained}

在这个示例中，我们要找一条内接于一个三角形的抛物线，即与三角形三边都相切的抛物线。
我不打算细讲得到该抛物线的第一部分。你会注意到这一行

\begin{mybox}
\begin{verbatim}
 L.euler = T:euler_line():swap_line()
\end{verbatim}
\end{mybox}

它交换了欧拉线的两个端点，稍后我们会看到这一点。

为了构造切点，需要求出抛物线与各边的交点：

\begin{mybox}
\begin{verbatim}
 z.ta = intersection(PA, T.bc)
 z.tb = intersection(PA, T.ca)
 z.tc = intersection(PA, T.ab)
\end{verbatim}
\end{mybox}

下面我们详细说明如何求直线 $(ab)$ 与抛物线的交点。这里欧拉线充当抛物线的准线。它的两个端点已被交换，以保持横坐标的正确顺序——即左侧为负、右侧为正。

为简化计算，把坐标系平移到以抛物线顶点为原点是有益的。从 $K$ 指向 $F$ 的焦轴（\code{major\_axis}）成为纵轴，而横轴的选取使新坐标系为 \code{右手系}。

我在代码中保留了 |z.U = OCCS.x| 和 |z.V = OCCS.y|，以便可视化新坐标系，例如用 |\tkzDrawSegments[red,->](S,U S,V)|。这个新坐标系是这样创建的：

\begin{mybox}
\begin{verbatim}
 O.SKF  = occs(L.KF,z.S)
\end{verbatim}
\end{mybox}

直线 $(KF)$，即抛物线的对称轴，成为纵轴。在这个新坐标系中，抛物线的方程为 $y = \dfrac{x^2}{2p}$，其中 $p$ 是距离 $KF$，也称为\emph{通径}。

\code{occs} 类的 \code{coordinates} 方法可用来取得每个点在新坐标系中的坐标。\code{param\_line} 函数计算直线方程的系数（该函数目前是内部的，其名字可能改变）。然后用 \code{solve\_para\_line} 求直线与抛物线的公共点（同样，该函数是内部的，可能修改）。

得到的结果是两个横坐标，必须放在过 $S$ 且与焦轴正交的那条轴上。这就是为什么正确放置曲线很重要。如果你对欧拉线去掉 |swap_line|，会看到曲线变成了原来的镜像。抛物线整体不变，但交点会变。

最后，放置这些交点的横坐标，然后求过这些横坐标、与欧拉线正交的直线与直线 $(ab)$ 的交点。

注意：这种几何方法比用公式直接求交点坐标更合适。因为那些坐标会表示在新坐标系中，还需要再做一次变换才能回到原坐标系。

\begin{verbatim}
\directlua{
   init_elements()
   z.A = point(0, 0)
   z.B = point(6, 0)
   z.C = point(2, 3)
   T.ABC  = triangle(z.A, z.B, z.C)
   L.euler = T.ABC:euler_line():swap_line()
   z.F = T.ABC:kimberling(110)
   z.H = T.ABC.orthocenter
   z.O = T.ABC.circumcenter
   z.Ω = point(0, 0)
   z.i = point(1, 0)
   z.j = point(0, 1)
   CO.PA  = conic(z.F, L.euler, 1)
   PA.curve = CO.PA:points(-3.5, 5.5, 50)
   local p = CO.PA.p
   z.K  = CO.PA.K
   z.S  = tkz.midpoint(z.F,z.K)
   L.KF = CO.PA.major_axis
   z.ta = intersection(CO.PA,T.ABC.bc)
   z.tb = intersection(CO.PA,T.ABC.ca)
   z.tc = intersection(CO.PA,T.ABC.ab)
% new occs
  O.SKF = occs(L.KF, z.S)
  z.U = O.SKF.x
  z.V = O.SKF.y
% line a,b
  z.a = point(3, 6)
  z.b = point(8, -1)
  L.ab = line(z.a, z.b)
% coordinates in the new occs
  Xa,Ya = O.SKF:coordinates(z.a)
  Xb,Yb = O.SKF:coordinates(z.b)
% solve in the new occs
  local r,s = tkz.line_coefficients(Xa, Ya ,Xb, Yb)
  r1,r2 = tkz.solve_quadratic_(1, -2 * p * r, -2 * p * s)
  z.x = O.SKF.abscissa:report(r1, z.K)
  z.y = O.SKF.abscissa:report(r2, z.K)
  L1 = L.euler:ortho_from(z.x)
  L2 = L.euler:ortho_from(z.y)
  z.s_1 = intersection(L.ab, L1)
  z.s_2 = intersection(L.ab, L2)}
  \begin{tikzpicture}
  \tkzGetNodes
   \tkzDrawCoordinates[smooth,purple,thick](PA.curve)
   \tkzDrawLines[add = .2 and 1](A,B A,C B,C K,F O,H)
   \tkzDrawPolygon[thick,cyan](A,B,C)
   \tkzDrawSegment[blue](a,b)
   \tkzDrawPoints(F,K,H,S,O)
   \tkzDrawPoints(A,B,F,K,S,ta,tb,tc)
   \tkzDrawPoints[red,size=2](s_1,s_2)
   \tkzLabelPoints[red,above](s_1)
   \tkzLabelPoints[red,right](s_2)
   \tkzLabelPoints(F,S,O,A,B)
   \tkzLabelPoints[above](C)
   \tkzLabelPoints[left](H,K)
  \end{tikzpicture}
\end{verbatim}

\directlua{
   init_elements()
   z.A = point(0, 0)
   z.B = point(6, 0)
   z.C = point(2, 3)
   T.ABC  = triangle(z.A, z.B, z.C)
   L.euler = T.ABC:euler_line():swap_line()
   z.F = T.ABC:kimberling(110)
   z.H = T.ABC.orthocenter
   z.O = T.ABC.circumcenter
   z.Ω = point(0, 0)
   z.i = point(1, 0)
   z.j = point(0, 1)
   CO.PA  = conic(z.F, L.euler, 1)
   PA.curve = CO.PA:points(-3.5, 5.5, 50)
   local p = CO.PA.p
   z.K  = CO.PA.K
   z.S  = tkz.midpoint(z.F,z.K)
   L.KF = CO.PA.major_axis
   z.ta = intersection(CO.PA,T.ABC.bc)
   z.tb = intersection(CO.PA,T.ABC.ca)
   z.tc = intersection(CO.PA,T.ABC.ab)
% new occs
  O.SKF = occs(L.KF, z.S)
  z.U = O.SKF.x
  z.V = O.SKF.y
% line a,b
  z.a = point(3, 6)
  z.b = point(8, -1)
  L.ab = line(z.a, z.b)
% coordinates in the new occs
  Xa,Ya = O.SKF:coordinates(z.a)
  Xb,Yb = O.SKF:coordinates(z.b)
% solve in the new occs
  local r,s = tkz.line_coefficients(Xa, Ya ,Xb, Yb)
  r1,r2 = tkz.solve_quadratic_(1, -2 * p * r, -2 * p * s)
  z.x = O.SKF.abscissa:report(r1, z.K)
  z.y = O.SKF.abscissa:report(r2, z.K)
  L1 = L.euler:ortho_from(z.x)
  L2 = L.euler:ortho_from(z.y)
  z.s_1 = intersection(L.ab, L1)
  z.s_2 = intersection(L.ab, L2)}

\begin{center}
\begin{tikzpicture}
  \tkzGetNodes
  \tkzDrawCoordinates[smooth,purple,thick](PA.curve)
  \tkzDrawLines[add = .2 and 1](A,B A,C B,C K,F O,H)
  \tkzDrawPolygon[thick,cyan](A,B,C)
  \tkzDrawSegment[blue](a,b)
  \tkzDrawPoints(F,K,H,S,O)
  \tkzDrawPoints(A,B,F,K,S,ta,tb,tc)
  \tkzDrawPoints[red,size=2](s_1,s_2)
  \tkzLabelPoints[red,above](s_1)
  \tkzLabelPoints[red,right](s_2)
  \tkzLabelPoints(F,S,O,A,B)
  \tkzLabelPoints[above](C)
  \tkzLabelPoints[left](H,K)
\end{tikzpicture}
\end{center}
\endinput